\paragraph*{04.03.02.04 Das Projekt mit den Supportfunktionen in Einklang bringen}
\kcilabel{para:04.03.02.04_Supportfunktionen}{04.03.02.04}
\label{para:04.03.02.04_Supportfunktionen}

%Task
\par Aufgrund seines Volumens und der Beteiligung von sieben externen Partnern stellte das Programm \gls{r2h} hohe Anforderungen an die Supportfunktionen des \gls{adac} und der \gls{ais}. Meine Aufgabe war es, eine nahtlose Zusammenarbeit mit der \gls{zek}, der \gls{jur} und der \gls{con} sicherzustellen, um die operativen Lieferbedarfe des \gls{ws} \gls{ricefw} rechtlich und betriebswirtschaftlich abzusichern und Verzögerungen durch administrative Flaschenhälse zu vermeiden.

%Action
\par Ich initiierte einen engen Abstimmungsprozess mit dem \gls{zek}, um Leistungsbeschreibungen und Ausschreibungen nach meinen Governance-Standards zu präzisieren. Mit der Rechtsabteilung sicherte ich Datenschutz, \gls{gob} und Datensicherheit bei Anbindung externer Teams. Mit dem Projekt-Controlling etablierte ich Budget-Monitoring für die \gls{ricefw}-Umsetzung. Implementierte Prozess für zeitnahe Leistungsbestätigung, um Rechnungsprüfung auf Basis qualitätsgeprüfter Deliverables zu ermöglichen.

%Result
\par Administrative Prozesse wurden synchron zum Projektfortschritt gehalten; keine Verzögerungen bei Beauftragung oder Bezahlung, was Motivation und Liefertreue stärkte. Compliance gegenüber konzernweiten Vorgaben vollumfänglich gewahrt durch Einbindung von \gls{jur}, \gls{zek}, \gls{con}, \gls{jdr}, \gls{rec} und \gls{com}. Damit habe ich nachgewiesen, dass ich komplexe IT-Projekte in administrative Rahmenbedingungen integrieren und diese als Enabler nutzen kann.